% !TeX program = xelatex
% Compile with XeLaTeX (twice for the table of contents).
\documentclass[12pt,a4paper]{article}
\usepackage[margin=1in]{geometry}
\usepackage{fontspec}
\usepackage[parfill]{parskip}
\usepackage{xeCJK}
\IfFontExistsTF{Noto Serif CJK SC}{\setCJKmainfont{Noto Serif CJK SC}}{\IfFontExistsTF{SimSun}{\setCJKmainfont{SimSun}}{\setCJKmainfont{FandolSong-Regular}}}
\usepackage{hyperref}
\setlength{\parindent}{0pt}
\title{Putting the TCM in NUS: A Traditional Chinese Medicine Dietary Formulation for Insomnia, Fatigue, Stress, and Loss of Appetite in University Students}
\author{Ryan Chen Rui Yang\\GEC1044 Final Report}
\date{7 November 2026}
\begin{document}
\maketitle

\begin{abstract}
University students in tertiary institutions frequently experience a common set of symptoms: physical and cognitive fatigue, sleep architecture disruption (insomnia), chronic affective tension (stress), and loss of appetite (anorexia). This report uses traditional Chinese medicine (TCM) to consider how these complaints might be understood together. It proposes a congee containing lily bulb, wolfberry, red date, Chinese angelica, lotus seed and cultivated cordyceps, with rice, ginger and water. The ingredients are analysed through their traditional properties, their proposed roles in the formula and their practical relevance to student life. Course readings on food as medicine, nourishing life and the organisation of prescriptions provide the historical and conceptual basis. The resulting congee brings traditional food practices into the context of student life, with attention to practical use and individual suitability.
\end{abstract}
\pagebreak

\tableofcontents
\clearpage

\section{Introduction}
Nowadays, students find themselves having to balance their academics, social life, and their own rest. This leads to a general lack of sleep, restlessness, and possibly a loss of appetite as well. However, not all students will experience these same symptoms, and each individual's personal issues can have a range of underlying causes.

From a TCM perspective, it is possible to look into how such experiences can possibly relate to one another. TCM practitioners tend to consider a trend in circumstances and changes over time to an individual to assess their health. As such, this report uses one possible pattern, Heart and Spleen deficiency with Liver qi constraint, as a framework for designing a food-based formula. This pattern provides a starting point for exploring the recipe's design.

But why are the particular ingredients in a meal for a student with this combination of complaints together? To answer this question, it is not enough to state the properties of each ingredient. The properties must also be connected to the primary purpose of the formula (the principle behind the formula itself), what tensions could come about when they are combined, and which benefits are still under speculation.

\newpage
\section{Understanding the Student's Complaints through TCM}
\subsection{The Body as an Interconnected System}
In classical Chinese medicine, the body is seen through the lens of the relationships among qi, yin and yang, as well as the organ systems, known as zang and fu. These organ systems refer to various functions and connections in the body, and do not necessarily represent modern anatomical organs or their functions. It is indicated that the Suwen compares the organs to officials in a state: the Heart is associated with the ruler, while other organs have their own responsibilities. According to this analogy, an individual's health is more so determined by the coordinated activity of the entire system, rather than just one part alone.\footnote{Week 2 lecture video, Organs and Meridians in Classic Chinese Medical Texts.}

In TCM, it states that the Spleen has a responsibility of transforming food and transporting the nutrients throughout the body, while the Liver ensures the smooth flow of qi. The Spleen's connection to blood and the shen (also often referred to as the spirit or mind) is through the Heart. These TCM concepts crucially concern the organs but in a much more broader sense. In other words, they refer to the roles and relationships of the organs in TCM rather than their structures, as is the case in Western medicine.

\subsection{Study, Emotion and Daily Habits}
The seven emotions are discussed in some of the lecture material, including the idea of pensiveness. It was shown how too much of an emotion, or having an emotion for a long period of time can contribute to illness. Specific readings also identified how physical fatigue could be caused by irregular eating and sleeping patterns.\footnote{Week 1 lecture video, Beginnings of Medicine; Week 10 reading, Nourishing Life, Cultivation and Material Culture.} From a TCM perspective, excessive worry and long study periods can cause Spleen qi constraint. This can be used to explain why poor appetites and fatigue are often linked. If the Spleen's transforming function is compromised, this can cause potential issues with the body's digestion and obtainment of necessary nutrients.

Liver qi constraints can also be of concern. When students get frustrated and feel pressured over their current struggles, they can end up feeling irritable and tense and create digestion issues. In TCM, this can be explained by the Wood-Earth relationship: if the Liver qi is constrained, it can affect the function of the Spleen. This indicates that emotional stress and digestive issues are linked, but the extents of its effects can vary between students.

Poor sleep can also play a part in this. In TCM, the Heart is said to house our Shen, and having enough blood in the Heart is said to keep the mind relaxed. Thus, Heart and Spleen Deficiency can be used to explain why students can suffer from fatigue and poor appetites.

This gives a baseline as to why the recipe has been suggested. However, note that this cannot be applied to all students who suffer with these symptoms. A TCM professional would need to look at their lifestyle, emotional status and other symptoms to make a proper judgement.

\newpage
\section{Recipe Design and Justification}
\subsection{Food, Medicine and Nourishing Life}
The course reading by Ping-Chung Leung describes a Chinese medical approach in which food and selected herbs are used to support health in relation to a person's condition. Its list places lotus seed and lily bulb among ingredients used for general support, angelica and wolfberry among those used to support blood, and cordyceps among ingredients used to supplement yin.\footnote{Week 10 reading, Ping-Chung Leung, Health, Wellbeing, Competence and Aging, chapter 1, table 1.1.} These groupings provide a useful starting point for considering how the ingredients might complement one another.

The idea of yangsheng, or nourishing life, gives the meal a wider context. One course reading describes a tradition that included diet, exercise and attention to daily and seasonal rhythms.\footnote{Week 2 reading, Michael Stanley-Baker, “Daoism and Medicine.”} For a student, this encourages attention to regular meals and rest as connected habits. The congee fits within this broader approach alongside rest, regular meals and attention to the pressures of student life.

\subsection{Proposed Recipe for Congee}
The following original formulation brings all the above ideas together in a meal specially designed for students, meant to be simple to prepare, easy to eat, and require little preperation time to accomodate for their busy schedules.

\begin{itemize}
\item Dried lily bulb (bai he), 15 g.
\item Dried wolfberry fruit (gouqizi), 15 g.
\item Red dates (da zao), five, pitted and sliced.
\item Angelica root (dang gui), 3 g.
\item Lotus seeds (lian zi), 15 g, with green cores removed.
\item Cultivated \emph{Cordyceps militaris} fruiting body, 2 g.
\item Rice, 80 g; fresh ginger, one thin slice; and water, approximately 1 litre.
\item Optional: a little sugar or honey for sweetness.
\end{itemize}

1) Rinse the rice, lily bulb and lotus seeds, then soak the dried lily bulb and lotus seeds until softened. 

2) Bring the water to a boil, and then add the rice, lily bulb, lotus seeds, red dates, angelica, cordyceps and ginger and simmer for approximately 45--50 minutes or till the ingredients are soft. 

3) In the last 5 minutes, add in the wolfberries. The congee can be served immediately after, and sugar or honey can be added for a sweeter taste.

Note: Congee is generally a go-to option when one is feeling unwell or has a loss of appetite.

\subsection{Ingredient Analysis}
\subsubsection*{Lotus seed and lily bulb: linking appetite and sleep}
Lotus seed has traditionally been described as a sweet and astringent ingredient, which generally targets the Spleen, Heart and Kidney. In our recipe, it addresses two concerns: supporting the Spleen relates to appetite, while nourishing the Heart and settling shen relates to sleep. The Week 10 reading on health and ageing also lists lotus seed (Semen Nelumbinis) among foods and herbs used for general support.\footnote{Week 10 reading, Health, Wellbeing, Competence and Aging, chapter 1, table 1.1.} This makes it a useful foundation for a nourishing meal, alongside its more specific role in the proposed pattern. However, the fact that it is an astringent calls for caution in a person that is experiencing constipation.

Next, lily bulbs, also known as Baihe are traditionally described as a slightly cold and sweet ingredient, and primarily targets the Heart and Lung, meant to nourish yin and calm restlessness. The same Week 10 reading lists lily bulb (Bulbus Lilii) alongside lotus seed for general support. Another Week 10 reading (pursuing the elixir of life) mentions Baihe as a food with calming effects in its discussion of yangxin, or cultivating a calm Heart.\footnote{Week 10 readings, Health, Wellbeing, Competence and Aging, chapter 1, table 1.1; Longevity and Vitality, discussion of yangxin.} 

This link helps explain why lily bulb suits a meal intended for a student whose stress and sleep difficulties occur together. When paired with lotus seed, it connects nourishment with the aim of settling the Shen at night. 

\subsubsection*{Wolfberry and angelica: support for blood and yin}
Wolfberry (goji berry) is traditionally known as a sweet and neutral fruit, and is generally associated with the Liver and Kidney. Its proposed role is to nourish Liver and Kidney yin and support blood. A course reading discusses the addition of wolfberry and chrysanthemum to the six-ingredient Rehmannia pill for a pattern involving Liver-yin weakness and blurred vision.\footnote{Week 7 reading, Principles of Chinese Medicine, discussion of Liuwei Dihuang Wan and Qiju Dihuang Wan.} This example shows how an existing formula can be adapted for a particular pattern. It also offers a useful point of comparison for a student who spends long periods reading: both cases draw attention to the eyes within a wider account of nourishment.

Angelica herb (dang gui) is also an ingredient kown to be traditionally warm, sweet and pungent, and is associated with the Liver, Heart and Spleen. It is included to nourish and move blood, linking the account of Heart blood and sleep with the Liver's role in movement. Leung's reading groups angelica with ingredients used to support blood. This classification helps explain its supporting role in a formula centred on the Heart and Spleen.\footnote{National Center for Complementary and Integrative Health, “Menopausal Symptoms: In Depth,” \url{https://www.nccih.nih.gov/health/menopausal-symptoms-in-depth}, accessed 7 November 2026.}

Note: It is highly advised for angelica herb and blood-thinning medication to not be taken together, as it poses a dangerous risk to the individual.

\subsubsection*{Red dates and ginger: making the formula workable}
Traditionally a sweet and warming herb, red date is added for its effect on the Spleen and Stomach. It will 'centralise' the dish and help balance all the ingredients inside it. Its sweet nature also helps make the meal more pleasant to eat, which can helps individuals experiencing poor appetite. Note that the sweetness should be moderated for those that are managing blood glucose levels.

Fresh ginger is also known to be a warm and pungent ingredient that is associated with the Spleen, Stomach and Lung. Its minute presence in the dish helps to balance out the richer ingredients. It was also mentioned in a week 9 reading that ginger was used in a southern Chinese postpartum broth as an ingredient to counter cold.\footnote{Week 9 reading, The Chinese Way of Health through Food, discussion of postpartum broth.} In our congee, the ginger serves a similar warming purpose.

\subsubsection*{Cultivated cordyceps: a qualified choice for fatigue}
Cordyceps in TCM is related to the Lung and Kidney meridians and can therefore be used as an adjunctive ingredient to help combat fatigue. In Leung's reading, he lists cordyceps among ingredients that are used to supplement yin.\footnote{Week 10 reading, Health, Wellbeing, Competence and Aging, chapter 1, table 1.1.} This ingredient is specifically a cultivated form of \emph{Cordyceps militaris} fruiting body, and should not be confused with mycelium or the traditional caterpillar fungus, \emph{Ophiocordyceps sinensis}.

For our specific target audience, it is important to ask whether it is truly worth paying the extra cost of using this ingredient, due to its traditional use.\footnote{Memorial Sloan Kettering Cancer Center, “Cordyceps,” \url{https://www.mskcc.org/cancer-care/integrative-medicine/herbs/cordyceps}, accessed 7 November 2026. The source discusses cordyceps products and interactions; evidence for a specific cultivated product may differ.}

Note: Like angelica herb, cordyceps products can interact with anti-coagulants or anti-diabetic drugs, so it is advisable to seek medical advice before use.

\subsubsection*{Rice and water: the core of congee}
Rice and water form the base of the congee, making the herbs part of a warm and foundational meal that is easy to prepare and eat.
\subsection{How the Ingredients Work Together}
The chinese formula design often constructs ingredients using a hierarchy approach: monarch ( jun ), minister ( chen ), assistant ( zuo ) and guide ( shi ). This was seen in the 6 ingredient Rehmannia pill example in the course reading, where 1 ingredient can be used as a leader whilst the others help or counteract any unwanted negative effects. It also notes that not all formulas will have all 4 hierarchies.\footnote{Week 7 reading, Principles of Chinese Medicine, discussion of Liuwei Dihuang Wan.}

In our case, lotus seed and lily bulb are proposed as the leading pair because they link the Spleen and Heart concerns central to the case. Wolfberry, angelica and cordyceps are the supporting herbs that help with the blood, yin and fatigue. Dates and ginger help to add flavour, while rice and water have been used as the base. These roles express the reasoning behind this particular recipe and show how its ingredients support a shared purpose.

The relationship also has a tension. Rich, sweet ingredients may be appealing and fit a tonifying approach, yet a student who already feels full after eating may not tolerate a large or rich bowl. Ginger is included for balance within TCM theory, while portion size and individual preference also matter. Sometimes a lesser dose, or less complex dish will be more suitable.Likewise, adding more herbs does not necessarily make the formula stronger, and you will have to consider if it is worth paying extra, longer preparation time and possible interaction with other ingredients.
\newpage
\section{Discussion and Critical Reflection}
\subsection{Prevention and Its Limits}
The Suwen introduces the idea of managing sickness before it actually develops, and the course reading frames this idea within a broader culture of sustaining health for the long-term.\footnote{Week 2 reading, Michael Stanley-Baker, “Daoism and Medicine,” discussion of Suwen 2.12 and yangsheng.} For our target group of students, this can be a form of encouragement for them to not only pay attention to their academics, but also to how they manage their sleep, meals and stress levels, to prevent any further deprovement. As such, the congee recipe is only one small part of a broader approach to improving daily habits.

Rather than viewing this recipe as a quick fix for fatigue, it is far more effective to consider it as part of a structured daily routine to better improve one's health. Combining this meal with a regular sleep schedule yields much better health outcomes than using it as a temporary fix to justify staying up late to study.


\subsection{Traditional Rationale and Clinical Evidence}
In TCM, the terminologies explain why certain ingredients are selected based on traditional concepts, while modern clinical research evaluates the effects of ingredients through various different scientific methods. As a result, some terms that are identical in both systems could have different meanings. For example, "dampness" does not literally mean that a person is wet or soggy, as it refers to the presence of heavy sluggish energy.

The way these terms are used in TCM can be a bit of a culture shock if their meanings are taken too literally, and thus it is important to make distinctions between both medical systems.

Another distinction is that modern healthcare represents sleep, diet, social aspects, and mental wellbeing into a holistic framework, while TCM represents these elements into its own network of interconnected organ systems and patterns. Throughout GEC1044, it was shown how different medical systems organise and structure their clinical observations towards particular treatment choices.

The proposed recipe brings together a bunch of ingredients, each with their own set of traditional benefits. Like many other types of research, more can be explored about how this particular combination relates to student sleep and wellbeing.

\subsection{Individual Suitability and Necessary Precautions}
This recipe may not be suitable for every person that experiences tiredness, bad sleep quality, and loss of appetite. If unsure, always seek professional advice before attempting to prepare and consume such a recipe. As herbal ingredients are also known to have potential interactions with western medicines, each individual's health history and current medications need to be considered on a case-by-case basis.

Persistent, worsening or severe symptoms warrant assessment by a qualified healthcare professional. Urgent symptoms such as gastrointestinal bleeding or thoughts of self-harm require urgent medical help. The meal is not a substitute for diagnosis or treatment. This boundary is consistent with the report's central argument: attending to the whole person includes recognising when food alone is insufficient.

\newpage
\section{Conclusion}
The proposed congee connects four student complaints through one possible TCM pattern: difficulty with digestion and nourishment, unsettled sleep, and constrained movement of qi. Lotus seed and lily bulb form the leading pair; wolfberry, angelica and cordyceps extend the traditional rationale; and dates, ginger, rice and water make the proposal a meal. Their roles become meaningful through their relationships to the chosen pattern and to one another.

The course readings show that food can be considered alongside medicinal practice and daily cultivation, while historical examples illustrate how ingredients and formulas are adapted to particular circumstances. Together, these ideas show how an everyday meal can be designed around a connected understanding of sleep, appetite and fatigue, while remaining attentive to individual needs, cost and safety.

\newpage
\section{Bibliography}
GEC1044. Lectures. Course notes supplied in Lectures.docx.

GEC1044. Readings. Course reading compilation supplied in Readings.docx.

Leung, Ping-Chung. Health, Wellbeing, Competence and Aging. Singapore: World Scientific, 2016. Chapter 1 reproduced in GEC1044, Readings.

Memorial Sloan Kettering Cancer Center. “Cordyceps.” \url{https://www.mskcc.org/cancer-care/integrative-medicine/herbs/cordyceps}. Accessed 7 November 2026.

National Center for Complementary and Integrative Health. “Menopausal Symptoms: In Depth.” \url{https://www.nccih.nih.gov/health/menopausal-symptoms-in-depth}. Accessed 7 November 2026.

National Center for Complementary and Integrative Health. “Sleep Disorders and Complementary Health Approaches: Usefulness and Safety.” \url{https://www.nccih.nih.gov/health/sleep-disorders-and-complementary-health-approaches}. Accessed 7 November 2026.

Stanley-Baker, Michael. “Daoism and Medicine.” In Routledge Handbook of Chinese Medicine, edited by Vivienne Lo, Michael Stanley-Baker and Dolly Yang, 173--187. London: Routledge, 2022. Reproduced in GEC1044, Readings.
\end{document}
